% !TEX root = ../main.tex
\chapter{マイグレーション証明チェックリスト}
\label{app:migration-proof-checklist}

\begin{chapterabstract}
本付録は、データ構造、API レスポンス、設定ファイル、DB スキーマなどを変更するときに、Lean 4 で何を証明すべきかを決めるための実務チェックリストです。
本文の各章では、同型、仕様弱化、バッチ移行、自然性、DB クエリ保存を個別に扱いました。
本付録では、それらをレビューで使える順序に並べ直します。
目的は、マイグレーションを「移せそうか」ではなく、「どの型からどの型へ変換し、どの性質が保存されるか」として記録できるようにすることです。
\end{chapterabstract}

\begin{learninggoals}
\item 変更前後の型、変換関数、逆変換、前提条件を分けて記録できるようになります。
\item 完全同型、片方向互換、制約付き同型、仕様弱化、情報損失がある移行を区別できるようになります。
\item Lean で証明すべき性質を、round-trip、バッチ保存、自然性、クエリ保存、実装差分の観点から選べるようになります。
\item 証明モデルが実装コードの何を保証し、何を保証していないかを、レビュー資料に残せるようになります。
\end{learninggoals}

\section{この付録を使う場面}

\begin{intuitionbox}
マイグレーション証明で最初に確認するのは、証明テクニックではありません。
最初に確認するのは、「何を何へ変換しているのか」と「何が保存されれば安全といえるのか」です。
型、関数、前提条件、保存したい性質を分けて書くと、証明すべき対象が小さくなります。
\end{intuitionbox}

本付録のチェックリストは、次のような変更で使います。

\begin{itemize}
\item ユーザー、注文、請求、ログなどのレコード型を変更します。
\item API レスポンスの wrapper、エラー形式、metadata の持ち方を変更します。
\item DB テーブルを分割、統合、正規化、非正規化します。
\item 設定ファイル、JSON、YAML、イベントスキーマをバージョンアップします。
\item 旧データを新形式で読み込みますが、完全な rollback を保証できるかは不明です。
\item バッチ移行後に件数、キー、集計値、join 結果が保存されるかを確認します。
\end{itemize}

\begin{softwarebox}
この付録は、実装前の設計レビュー、移行 PR のレビュー、CI に入れる Lean 定理の選定、移行後の監査資料に使えます。
すべての本番コードを Lean へ移す必要はありません。
代わりに、移行仕様の核だけを小さな検証モデルとして切り出します。
\end{softwarebox}

\FigurePlaceholder{fig-appD-migration-proof-flow.png}{0.86}{マイグレーション証明フロー}{fig:appD-migration-proof-flow}

図\ref{fig:appD-migration-proof-flow}は、この付録の読み方を表しています。
まず変更前後の型を決め、変換関数と逆変換の有無を確認します。
次に、情報損失と well-formed 条件を評価し、最後にバッチ移行、業務ロジック、実装差分へ確認範囲を広げます。

\section{最初に書く一枚メモ}

マイグレーション証明を始める前に、次の一枚メモを作ります。
このメモを作れない場合は、まだ Lean で証明する段階ではありません。
証明する前に、移行仕様そのものを整理する必要があります。

\begin{definitionbox}
\textbf{マイグレーション証明の一枚メモ}：次の五つを同じ場所に書いたものです。

\begin{enumerate}
\item 変更前の型 \(Old\)
\item 変更後の型 \(New\)
\item 旧から新への変換 \(migrate : Old \to New\)
\item 新から旧への変換、または部分的な逆変換 \(rollback : New \to Old\)
\item 保存したい性質、または弱めた仕様
\end{enumerate}
\end{definitionbox}

\begin{longtable}{@{}p{0.17\linewidth}p{0.32\linewidth}p{0.43\linewidth}@{}}
\toprule
項目 & 書くこと & 例 \\
\midrule
\endhead
変更前の型 & 旧スキーマ、旧DTO、旧レスポンス、旧テーブル行 & \texttt{UserV1}, \texttt{OldResponse A}, \texttt{OrderRowV1} \\
変更後の型 & 新スキーマ、新DTO、新レスポンス、新テーブル行 & \texttt{UserV2}, \texttt{NewResponse A}, \texttt{OrderRowV2} \\
正方向の変換 & 旧形式を新形式へ移す関数 & \texttt{migrate : UserV1 -> UserV2} \\
逆方向の変換 & 新形式を旧形式へ戻す関数であり、存在しない場合はその理由を書きます & \texttt{rollback : UserV2 -> UserV1} \\
前提条件 & 正常データ、空でない文字列、キー一意性、外部制約など & \texttt{WellFormedNew v} \\
主張 & 完全復元、片方向復元、正規化後一致、件数保存など & \texttt{rollback (migrate u) = u} \\
保証しないこと & Leanモデルの外にある制約、性能、DBエンジンの実挙動など & インデックス性能、外部API障害、並行更新 \\
\bottomrule
\end{longtable}

\section{チェックリスト本体}

この節は、レビューでそのまま使うためのチェックリストです。
各項目では、自然言語で答える欄と、Lean 上での形を対応させます。

\subsection{1. 変更前後の型は何か}

\begin{softwarebox}
型を決めるとは、単にクラス名やテーブル名を決めることではありません。
どのフィールドを持ち、どの値を許し、どの不変条件を持つデータとして扱うかを決めることです。
\end{softwarebox}

\begin{longtable}{@{}p{0.06\linewidth}p{0.32\linewidth}p{0.52\linewidth}@{}}
\toprule
確認 & 問い & 完了条件 \\
\midrule
\endhead
\(\square\) & 変更前の型を Lean で名前付きにしましたか & \texttt{structure UserV1 where ...} のように、旧形式が独立した型として書かれています。 \\
\(\square\) & 変更後の型を Lean で名前付きにしましたか & \texttt{structure UserV2 where ...} のように、新形式が独立した型として書かれています。 \\
\(\square\) & 単なる表現差と意味の差を分けましたか & フィールド名変更、ネスト化、単位変更、正規化、情報削除が区別されています。 \\
\(\square\) & 不変条件を型に含めるか、述語に分けるかを決めましたか & \texttt{WellFormedV1 : UserV1 -> Prop} のように、制約の置き場所が明示されています。 \\
\bottomrule
\end{longtable}

\begin{leanbox}
Lean では、まず本番の巨大な型をそのまま再現しません。
証明したい性質に必要なフィールドだけを持つ、小さな型を作ります。
この小さな型が、証明モデルです。
\end{leanbox}

\begin{listing}[htbp]
\begin{minted}{lean4}
structure UserV1 where
  id : Nat
  name : String

deriving Repr

structure Profile where
  displayName : String

deriving Repr

structure UserV2 where
  id : Nat
  profile : Profile

deriving Repr
\end{minted}
\caption{\texttt{UserV1}・\texttt{Profile}・\texttt{UserV2}}
\label{lst:appD_old_new_types}
\end{listing}

このコードでは、旧形式\texttt{UserV1}と新形式\texttt{UserV2}を、移行で比較したい最小の型として切り出しています。
\texttt{Profile}を別の構造体にしている点が、新形式では名前が入れ子のフィールドへ移ったことを表しています。

\subsection{2. 変換関数は何か}

変換関数は、移行仕様の中心です。
口頭で「移せる」と述べるだけでは不十分です。
Lean 上では、旧形式の値を受け取り、新形式の値を返す関数として書きます。

\begin{longtable}{@{}p{0.06\linewidth}p{0.32\linewidth}p{0.52\linewidth}@{}}
\toprule
確認 & 問い & 完了条件 \\
\midrule
\endhead
\(\square\) & 正方向の変換関数を書きましたか & \texttt{migrate : Old -> New} が定義されています。 \\
\(\square\) & 変換しないフィールドを明示しましたか & ID、キー、作成日時など、保存されるフィールドがコード上で確認できます。 \\
\(\square\) & 新しく作る値の根拠を説明しましたか & デフォルト値、派生値、外部参照値が区別されています。 \\
\(\square\) & 変換が失敗し得るかを決めましたか & 失敗しないなら \texttt{Old -> New}、失敗し得るなら \texttt{Old -> Option New} または \texttt{Old -> Except E New} を使います。 \\
\bottomrule
\end{longtable}

\begin{listing}[htbp]
\begin{minted}{lean4}
def migrateUser : UserV1 -> UserV2 :=
  fun u =>
    { id := u.id
      profile := { displayName := u.name } }
\end{minted}
\caption{\texttt{migrateUser}}
\label{lst:appD_migrate_function}
\end{listing}

このコードでは、旧形式の\texttt{id}をそのまま新形式の\texttt{id}へ移し、旧形式の\texttt{name}を\texttt{profile.displayName}へ入れています。
移行関数を読むときは、保存される値と、構造だけが変わる値とを分けて確認します。

\clearpage
\subsection{3. 逆変換は存在するか}

\begin{warningbox}
\texttt{Old -> New} と \texttt{New -> Old} の二つの関数があるだけでは、同型ではありません。
同型と述べるには、変換して戻したときに元に戻ることを証明する必要があります。
戻せない場合は、同型ではなく、片方向互換または仕様弱化として扱います。
\end{warningbox}

\begin{longtable}{@{}p{0.06\linewidth}p{0.32\linewidth}p{0.52\linewidth}@{}}
\toprule
確認 & 問い & 完了条件 \\
\midrule
\endhead
\(\square\) & 逆変換を書けますか & \texttt{rollback : New -> Old} を定義できます。 \\
\(\square\) & 逆変換を全入力について定義できますか & すべての \texttt{New} に対して戻せるか、well-formed な \texttt{New} に限るかを決めています。 \\
\(\square\) & 戻せない新形式の値がありますか & 新形式だけが持つ状態、任意の分割、外部生成 ID などが列挙されています。 \\
\(\square\) & 逆変換が失敗し得ますか & \texttt{New -> Option Old} にする必要があるかを判断しています。 \\
\bottomrule
\end{longtable}

\begin{listing}[htbp]
\begin{minted}{lean4}
def rollbackUser : UserV2 -> UserV1 :=
  fun u =>
    { id := u.id
      name := u.profile.displayName }
\end{minted}
\caption{\texttt{rollbackUser}}
\label{lst:appD_rollback_function}
\end{listing}

このコードでは、新形式で入れ子になった\texttt{profile.displayName}を、旧形式の\texttt{name}へ戻しています。
逆変換を書けることは出発点にすぎず、後続の round-trip 証明で本当に元へ戻るかを確認する必要があります。

\subsection{4. 情報損失はあるか}

情報損失がある移行を、完全同型として扱ってはいけません。
情報損失があるかどうかは、戻せるかどうかだけでなく、戻した値が業務上同じ意味を持つかどうかによっても判断します。

\begin{definitionbox}
本書では、\textbf{情報損失}を次のように読みます。
旧形式の値から新形式へ移した後、旧形式で区別できていた情報を復元できないなら、その移行には情報損失があります。
ただし、業務仕様としてその区別を不要にした場合は、完全復元ではなく、仕様弱化として証明対象を選び直します。
\end{definitionbox}

\begin{longtable}{@{}p{0.20\linewidth}p{0.34\linewidth}p{0.40\linewidth}@{}}
\toprule
移行の種類 & 典型例 & Leanで狙う性質 \\
\midrule
\endhead
表現変更だけ & フィールド名変更、ネスト化、順序の変更 & 両方向round-trip \\
旧から新へ拡張 & 新フィールドに固定デフォルトを入れる & 片方向round-trip、またはデフォルト条件付きround-trip \\
情報の分割 & \texttt{fullName} を \texttt{firstName} と \texttt{lastName} に分ける & well-formed条件付き復元、または表示文字列の互換性 \\
情報の削除 & 詳細ログから公開ログを作る & 完全復元ではなく、公開仕様の保存、単調性、安全な抽象化 \\
集約 & 複数明細から合計だけを残す & 合計保存、個別明細の復元不可を明示 \\
\bottomrule
\end{longtable}

\subsection{5. well-formed条件は必要か}

すべての値を許す型では証明できない場合でも、業務上正しい値に限れば証明できることがあります。
このとき必要になるのが、well-formed 条件です。

\begin{intuitionbox}
well-formed 条件は、「この形の値だけを本番データとして扱う」という約束です。
たとえば、空でない ID、重複しないキー、分割可能な氏名、対応する外部キーが存在する行などが該当します。
\end{intuitionbox}

\begin{longtable}{@{}p{0.06\linewidth}p{0.32\linewidth}p{0.52\linewidth}@{}}
\toprule
確認 & 問い & 完了条件 \\
\midrule
\endhead
\(\square\) & 型だけでは表せない制約がありますか & 空文字の禁止、キーの一意性、外部キーの整合性などが列挙されています。 \\
\(\square\) & 制約を述語として書きましたか & \texttt{WellFormedOld : Old -> Prop} や \texttt{WellFormedNew : New -> Prop} があります。 \\
\(\square\) & 移行が well-formed 性を保存しますか & \texttt{WellFormedOld x -> WellFormedNew (migrate x)} を証明対象にしています。 \\
\(\square\) & 条件付き復元の形を決めましたか & \texttt{h : WellFormedNew y} を仮定して復元性を主張します。 \\
\bottomrule
\end{longtable}

\begin{listing}[htbp]
\begin{minted}{lean4}
def WellFormedV1 (u : UserV1) : Prop :=
  u.name ≠ ""

def WellFormedV2 (u : UserV2) : Prop :=
  u.profile.displayName ≠ ""

theorem migrateUser_preserves_wellformed (u : UserV1)
    (h : WellFormedV1 u) :
    WellFormedV2 (migrateUser u) := by
  unfold WellFormedV1 at h
  unfold WellFormedV2 migrateUser
  exact h
\end{minted}
\caption{\texttt{WellFormedV1}・\texttt{WellFormedV2}・\texttt{migrateUser\_preserves\_wellformed}}
\label{lst:appD_wellformed}
\end{listing}

このコードでは、名前が空でないという業務上の条件を、旧形式と新形式それぞれの述語にしています。
\texttt{migrateUser\_preserves\_wellformed} は、旧形式でその条件を満たす値を移行すると、新形式でも入れ子の表示名が空でないことを示します。

\subsection{6. 証明したい性質はどの分類か}

証明したい性質を分類すると、Lean で書く定理の形が決まります。
分類しないまま始めると、完全同型を証明しようとして失敗し、その理由が仕様の問題なのか証明力の問題なのか分からなくなります。

\begin{longtable}{@{}p{0.18\linewidth}p{0.36\linewidth}p{0.40\linewidth}@{}}
\toprule
分類 & 意味 & 定理の形 \\
\midrule
\endhead
完全同型 & 新旧が同じ情報を持つ & \texttt{rollback (migrate x) = x} と \texttt{migrate (rollback y) = y} \\
片方向互換 & 旧データは新形式へ安全に移せる & \texttt{rollback (migrate x) = x} だけ、または旧仕様の保存 \\
制約付き同型 & well-formedな値の領域で相互に戻せる & 両方向の条件付きround-tripとwell-formed性の保存 \\
仕様弱化 & 完全復元ではなく、観測可能な性質だけを保存する & \texttt{observeNew (migrate x) = observeOld x} \\
安全な抽象化 & 詳細情報を公開可能な形へ落とす & \texttt{SafePublic (abstract x)}、単調性、非漏洩性 \\
\bottomrule
\end{longtable}

\begin{leanbox}
証明対象は、強いほどよいとは限りません。
実際には成り立たない完全同型を主張するよりも、業務上必要な観測結果だけを保存する定理のほうが正確です。
\end{leanbox}

\begin{listing}[htbp]
\begin{minted}{lean4}
theorem rollback_after_migrate (u : UserV1) :
    rollbackUser (migrateUser u) = u := by
  cases u
  rfl

theorem migrate_after_rollback (u : UserV2) :
    migrateUser (rollbackUser u) = u := by
  cases u with
  | mk id profile =>
    cases profile
    rfl
\end{minted}
\caption{\texttt{rollback\_after\_migrate}・\texttt{migrate\_after\_rollback}}
\label{lst:appD_roundtrip}
\end{listing}

このコードでは、旧形式から新形式へ移して戻す経路と、新形式から旧形式へ戻して移す経路の両方を証明しています。
二つの定理がそろって初めて、この例を完全同型として扱えます。

\subsection{7. バッチ化したとき性質は保存されるか}

1件の移行が安全でも、バッチ移行で必要な性質が自動的にすべて得られるわけではありません。
件数、順序、キー集合、重複の有無、エラー件数、集計値など、バッチ固有の性質を確認します。

\begin{longtable}{@{}p{0.06\linewidth}p{0.32\linewidth}p{0.52\linewidth}@{}}
\toprule
確認 & 問い & 完了条件 \\
\midrule
\endhead
\(\square\) & バッチ移行を \texttt{List.map migrate} で表せますか & 各要素を独立に変換するだけなら、\texttt{List.map} を使います。 \\
\(\square\) & 件数は保存されますか & \texttt{length (map migrate xs) = length xs} を証明するか、既存の定理で説明します。 \\
\(\square\) & 順序は保存されますか & \texttt{List.map} であれば順序保存が期待できます。
ソートやフィルタが入る場合は、別の仕様にします。 \\
\(\square\) & キーの一意性は保存されますか & ID 保存と旧キーの一意性から、新キーの一意性を導けるかを確認します。 \\
\(\square\) & 失敗する要素の扱いを決めましたか & 全体を失敗させるか、成功分だけを残すか、エラー一覧を返すかを決めます。 \\
\bottomrule
\end{longtable}

\begin{listing}[htbp]
\begin{minted}{lean4}
def migrateUsers (xs : List UserV1) : List UserV2 :=
  xs.map migrateUser

theorem migrateUsers_length (xs : List UserV1) :
    (migrateUsers xs).length = xs.length := by
  unfold migrateUsers
  simp
\end{minted}
\caption{\texttt{migrateUsers}・\texttt{migrateUsers\_length}}
\label{lst:appD_batch_migration}
\end{listing}

このコードでは、1件の移行関数 \texttt{migrateUser} を使います。
\texttt{List.map} は、その関数をリストのすべての要素へ適用します。
\texttt{migrateUsers\_}\allowbreak\texttt{length} は、移行によって件数が変わらないことを確認する、バッチ固有の性質です。

\subsection{8. 業務ロジックとの順序を入れ替えてもよいか}

API adapter やレスポンス wrapper の変更では、変換と業務ロジックの順序が重要になります。
先に移行してから業務ロジックを適用した場合と、先に業務ロジックを適用してから移行した場合とで結果が同じであれば、adapter は業務ロジックに干渉していないといえます。

\begin{definitionbox}
本付録で\textbf{順序を入れ替えてもよい}とは、二つの処理の合成結果が等しいことを意味します。
圏論では、可換図式や自然性として現れます。
Lean では、二つの式の等式として書きます。
\end{definitionbox}

\begin{longtable}{@{}p{0.06\linewidth}p{0.32\linewidth}p{0.52\linewidth}@{}}
\toprule
確認 & 問い & 完了条件 \\
\midrule
\endhead
\(\square\) & 変換前に適用する業務ロジックは何ですか & \texttt{businessOld : Old -> Old'} などの関数として名前が付いています。 \\
\(\square\) & 変換後に対応する業務ロジックは何ですか & \texttt{businessNew : New -> New'} などの関数として名前が付いています。 \\
\(\square\) & 二つの経路は同じ型へ到達しますか & 等式の左右の型が一致しています。 \\
\(\square\) & 順序交換は全入力について必要ですか & 全入力、成功ケースのみ、well-formed 入力のみのいずれかを決めています。 \\
\bottomrule
\end{longtable}

\clearpage

\begin{listing}[htbp]
\begin{minted}{lean4}
def renameV1 (u : UserV1) (s : String) : UserV1 :=
  { u with name := s }

def renameV2 (u : UserV2) (s : String) : UserV2 :=
  { u with profile := { displayName := s } }

theorem migrate_rename_commutes (u : UserV1) (s : String) :
    migrateUser (renameV1 u s) = renameV2 (migrateUser u) s := by
  cases u
  rfl
\end{minted}
\caption{\texttt{migrate\_rename\_commutes}}
\label{lst:appD_commute_business}
\end{listing}

このコードでは、名前変更を旧形式で行ってから移行する経路と、移行してから新形式で名前変更する経路とを比較しています。
\texttt{migrate\_rename\_commutes}は、二つの順序が同じ結果になることを示す可換性の例です。

\subsection{9. DBやクエリの結果は保存されるか}

DB スキーマ変更では、行の形を移せるだけでは不十分です。
本番で必要なクエリ結果が保存されるかを確認する必要があります。
ただし、実際の DB エンジン全体を Lean で完全に再現するのではなく、小さな有限モデルで必要な性質を表します。

\begin{longtable}{@{}p{0.06\linewidth}p{0.32\linewidth}p{0.52\linewidth}@{}}
\toprule
確認 & 問い & 完了条件 \\
\midrule
\endhead
\(\square\) & 保存したいクエリを選びましたか & すべての SQL ではなく、移行の安全性に関わる代表的なクエリを選びます。 \\
\(\square\) & クエリを関数としてモデル化しましたか & \texttt{queryOld : OldDB -> Result} と \texttt{queryNew : NewDB -> Result} を用意します。 \\
\(\square\) & 移行後の DB を作る関数を書きましたか & \texttt{migrateDB : OldDB -> NewDB} があります。 \\
\(\square\) & クエリ保存の等式を書きましたか & \texttt{queryNew (migrateDB db) = queryOld db} を定理として書きます。 \\
\(\square\) & モデル外の DB 仕様を明示しましたか & NULL、トランザクション、インデックス、照合順序などを保証外として記録します。 \\
\bottomrule
\end{longtable}

\begin{leanbox}
DB 移行の Lean モデルでは、テーブルをリスト、クエリを関数、整合条件を述語として表すことから始めます。
実際の DB の全機能を形式化するのではなく、保存したいクエリ結果だけを切り出します。
\end{leanbox}

\subsection{10. 証明モデルと実装コードの差分はどこか}

Lean で証明したことは、Lean で書いたモデルについての事実です。
したがって、証明モデルと本番実装がずれていれば、証明の実務的な意味は弱くなります。
証明した範囲と、証明していない範囲を必ず書きます。

\begin{warningbox}
「Lean で証明済み」と述べるときは、何が証明済みなのかを限定して書きます。
たとえば、「Lean モデル上では、\texttt{rollback (migrate u) = u} がすべての \texttt{UserV1} について成立します」と書きます。
「本番移行が完全に安全」とだけ書くと、外部 API、DB 制約、パーサ、並行実行、運用手順まで保証したように読めてしまいます。
\end{warningbox}

\begin{longtable}{@{}p{0.06\linewidth}p{0.32\linewidth}p{0.52\linewidth}@{}}
\toprule
確認 & 問い & 完了条件 \\
\midrule
\endhead
\(\square\) & Lean モデルに含めたフィールドを列挙しましたか & 証明対象のフィールドと、無視したフィールドが分かります。 \\
\(\square\) & 本番コードの変換処理と対応していますか & 実装関数、PR、設計書へのリンクを記録しています。 \\
\(\square\) & パース、シリアライズ、DB I/O を保証外にしましたか & 外部境界のテストが別に必要であると明記しています。 \\
\(\square\) & 証明が壊れたときの扱いを決めましたか & 仕様変更、モデル変更、実装バグのいずれかとしてレビューする方針があります。 \\
\(\square\) & CI に入れるかを決めましたか & \texttt{lake build} で検査するか、設計資料用にとどめるかが決まっています。 \\
\bottomrule
\end{longtable}

\section{証明タイプ別テンプレート}

この節では、レビュー時にそのまま記入できる形で、証明タイプ別のテンプレートをまとめます。

\subsection{完全同型テンプレート}

\begin{definitionbox}
完全同型として扱うのは、変更前後で情報が失われず、両方向の round-trip が成り立つ場合です。
\end{definitionbox}

\begin{listing}[htbp]
\begin{minted}{lean4}
-- Old, New を具体的な型に置き換える。
-- migrate : Old -> New
-- rollback : New -> Old

theorem rollback_migrate_id (x : Old) :
    rollback (migrate x) = x := by
  -- 型と定義に応じて cases, rfl, simp などを使う。
  sorry

theorem migrate_rollback_id (y : New) :
    migrate (rollback y) = y := by
  -- すべての New に対して戻せる場合だけ成立する。
  sorry
\end{minted}
\caption{\texttt{rollback\_migrate\_id}・\texttt{migrate\_rollback\_id}}
\label{lst:appD_template_iso}
\end{listing}

このテンプレート内の\texttt{sorry}は、読者が自分のプロジェクトで置き換える穴です。
本書本文では原則として\texttt{sorry}を残しませんが、ここではレビュー用テンプレートとして、証明を書く場所を明示しています。

\subsection{片方向互換テンプレート}

片方向互換では、旧データを新形式へ移しても旧仕様を再現できることを示します。
すべての新形式を旧形式へ戻せるとは限りません。

\begin{listing}[htbp]
\begin{minted}{lean4}
-- 旧データを移して戻すと、旧データに戻る。
theorem old_roundtrip (x : Old) :
    rollback (migrate x) = x := by
  sorry

-- または、旧仕様として観測できる値だけが保存される。
theorem old_observation_preserved (x : Old) :
    observeNew (migrate x) = observeOld x := by
  sorry
\end{minted}
\caption{\texttt{old\_roundtrip}・\texttt{old\_observation\_preserved}}
\label{lst:appD_template_one_way}
\end{listing}

このテンプレートでは、完全な二方向同型ではなく、旧データを移した後に旧形式向けの観測が保たれることを主張しています。
\texttt{old\_roundtrip}が強すぎる場合は、\texttt{old\_observation\_preserved}のように、保存したい観測へ弱めます。

\subsection{制約付き同型テンプレート}

制約付き同型では、すべての値ではなく、well-formed な値に限って復元性を主張します。

\begin{listing}[htbp]
\begin{minted}{lean4}
theorem migrate_rollback_on_wellformed
    (y : New) (h : WellFormedNew y) :
    migrate (rollback y) = y := by
  sorry

-- 二つの変換が制約領域の中に留まること。
theorem migrate_preserves_wellformed
    (x : Old) (h : WellFormedOld x) :
    WellFormedNew (migrate x) := by
  sorry
\end{minted}
\caption{\texttt{migrate\_rollback\_on\_wellformed}・\texttt{migrate\_preserves\_wellformed}}
\label{lst:appD_template_wellformed}
\end{listing}

このテンプレートの新形式側の定理に、前節の \texttt{old\_roundtrip} を組み合わせれば、両方向の復元性になります。
\texttt{migrate\_preserves\_wellformed} は、旧形式の制約領域から移した値が、新形式の制約領域に入ることを示します。
旧形式側も制約された部分だけを対象にするなら、\texttt{WellFormedNew y -> WellFormedOld (rollback y)} と、旧形式側の条件付き round-trip も追加する必要があります。
必要な両方向の復元性と制約保存がそろって初めて、制約された領域の同型として読めます。

\subsection{仕様弱化テンプレート}

完全復元が成り立たない場合は、業務上必要な観測だけを保存します。
たとえば、詳細ログを公開ログへ落とす場合、詳細ログそのものは戻せません。
しかし、公開してよい形式であることや、公開表示が期待どおりであることは証明できます。

\begin{listing}[htbp]
\begin{minted}{lean4}
-- normalizeOld と normalizeNew は、比較したい観測結果に落とす関数。
theorem normalized_observation_preserved (x : Old) :
    normalizeNew (migrate x) = normalizeOld x := by
  sorry

-- 情報を落とす変換なら、安全性述語を証明対象にする。
theorem abstraction_is_safe (x : Old) :
    SafePublic (abstract x) := by
  sorry
\end{minted}
\caption{\texttt{normalized\_observation\_preserved}・\texttt{abstraction\_is\_safe}}
\label{lst:appD_template_weakening}
\end{listing}

このテンプレートでは、元の値を完全に戻す代わりに、正規化後の観測や公開安全性だけを証明対象にしています。
\texttt{SafePublic} の定義が実際の公開方針を十分に表すかについては、別途レビューが必要であり、定理名だけで非漏洩性が得られるわけではありません。
情報を落とす移行では、このように「保存する性質」を明示してから Lean の定理へ落とします。

\section{レビュー用ワークシート}

次の表は、設計レビューまたは PR レビューに貼り付けるためのワークシートです。
チェック欄は、紙面や Markdown に転記しやすいように、単純な形にしています。

\begin{longtable}{@{}p{0.045\linewidth}p{0.20\linewidth}p{0.28\linewidth}p{0.31\linewidth}@{}}
\toprule
確認 & レビュー項目 & 記入内容 & Leanでの対応 \\
\midrule
\endhead
\(\square\) & 変更前の型 & 旧スキーマ名、主要フィールド、不変条件 & \texttt{Old}、\texttt{WellFormedOld} \\
\(\square\) & 変更後の型 & 新スキーマ名、主要フィールド、不変条件 & \texttt{New}、\texttt{WellFormedNew} \\
\(\square\) & 正方向変換 & 旧から新へ移す仕様 & \texttt{migrate : Old -> New} \\
\(\square\) & 逆方向変換 & 全域か、部分的か、存在しないか & \texttt{rollback : New -> Old} または \texttt{New -> Option Old} \\
\(\square\) & 情報損失 & 削除、集約、分割、外部生成値の有無 & 完全同型か仕様弱化かを選ぶ \\
\(\square\) & 前提条件 & 正常データ、キー制約、入力制約 & \texttt{WellFormed}述語 \\
\(\square\) & round-trip & どちら向きに戻せるか & \texttt{rollback (migrate x) = x} など \\
\(\square\) & バッチ性質 & 件数、順序、キー、集計値 & \texttt{List.map}、長さ保存、全要素性質 \\
\(\square\) & 業務ロジックとの可換性 & adapterの位置を変えてよいか & \texttt{migrate (businessOld x) = businessNew (migrate x)} \\
\(\square\) & DBクエリ保存 & 保存したい代表クエリ & \texttt{queryNew (migrateDB db) = queryOld db} \\
\(\square\) & 実装差分 & 証明モデルに含めない処理 & 保証外の明示、別テストへの接続 \\
\(\square\) & CI運用 & 証明をいつ検査するか & \texttt{lake build}、レビュー手順 \\
\bottomrule
\end{longtable}

\section{危険信号}

次のいずれかに該当する場合は、証明を書く前に仕様を見直します。

\begin{warningbox}
\begin{itemize}
\item 「互換性がある」という言葉だけで、どの性質が保存されるかが書かれていません。
\item 逆変換が存在すると述べていますが、戻せない新形式の値があります。
\item 情報を削除しているにもかかわらず、完全な round-trip を主張しています。
\item well-formed 条件が本番データの暗黙の約束になっており、Lean 上にも設計書にも書かれていません。
\item 1件の移行だけを確認し、バッチ時の件数、キー、失敗件数を確認していません。
\item adapter を入れる位置を変えていますが、業務ロジックとの可換性を確認していません。
\item Lean モデルのフィールドが、本番移行コードと対応していません。
\item 「証明済み」と書いていますが、証明した範囲と保証外の範囲が分かれていません。
\end{itemize}
\end{warningbox}

\section{証明済みと書くための文例}

最後に、設計書や PR に書ける文例を示します。
過度に強い表現を避け、Lean モデル上の保証として範囲を限定します。

\begin{softwarebox}
\textbf{完全同型の場合。}
Lean モデル上の\texttt{UserV1}と\texttt{UserV2}について、\texttt{rollbackUser (migrateUser u) = u} と \texttt{migrateUser (rollbackUser v) = v} を証明しました。
したがって、モデル化したフィールド \texttt{id} と \texttt{displayName} については、移行で情報を失いません。
パース、DB I/O、既存の本番データがすべて well-formed であることは、別途テストと移行前検査で確認します。
\end{softwarebox}

\begin{softwarebox}
\textbf{仕様弱化の場合。}
この移行は完全同型ではありません。
旧形式で保持していた詳細情報は、新形式から復元できません。
代わりに、業務上必要な観測値について、\texttt{normalizeNew (migrate x) = normalizeOld x} を Lean モデル上で証明しました。
そのため、本 PR で保証するのは完全復元ではなく、正規化後の表示互換性です。
\end{softwarebox}

\begin{softwarebox}
\textbf{バッチ移行の場合。}
1件の移行関数\texttt{migrateUser}を\texttt{List.map}でリスト全体へ持ち上げるモデルを作りました。
Lean モデル上では、移行前後でリスト長が保存されることを証明しました。
重複キー、外部キーの整合性、失敗レコードの扱いは、別の定理または移行前後の検査項目として管理します。
\end{softwarebox}

\section{本付録のまとめ}

マイグレーション証明は、変更前後の型、変換関数、逆変換、前提条件、保存したい性質を分けることから始まります。

完全同型、片方向互換、制約付き同型、仕様弱化、情報損失がある移行を混同しません。

Lean では、round-trip、well-formed 保存、バッチ保存、業務ロジックとの可換性、クエリ保存を小さな定理として表します。

証明モデルは、本番実装そのものではありません。
証明した範囲と保証外の範囲を、設計書や PR に明示します。

チェックリストは、証明を書くためだけでなく、移行仕様の曖昧さを見つけるためにも使えます。
